\section{Compiler Pipeline and Implementation Foundations}\label{sec:compiler-pipeline}

The Motivo compiler is organized as a strictly sequential multi-pass pipeline.
Each stage transforms the representation of the musical program, moving from abstract textual intent to a
concrete binary
artifact.
By decoupling these stages, the compiler catches errors early, keeps musical structure fully resolved before code
generation, and allows each module to be tested in isolation.

\subsection{Orchestration and Control Flow}\label{detail-subsec:orchestration}

The public entry point \texttt{motivo::compile} in \texttt{compiler/src/motivo/compiler.cpp} orchestrates the full
pipeline.
A single \texttt{DiagnosticsEngine} instance is created at the start and passed by reference into every stage that can
report structured errors; the engine accumulates \texttt{Diagnostic} records tagged with a \texttt{DiagnosticStage}
(\texttt{Parsing}, \texttt{Semantic}, \texttt{Lowering}, or \texttt{Output}) and an optional source location.
After each stage, \texttt{compiler.cpp} checks for stage-specific failures and returns early with the collected
diagnostics rather than proceeding on invalid input.

The orchestration sequence matches the implementation exactly:

\begin{enumerate}
  \item \texttt{parsing::parse\_stream} reads the input \texttt{FILE*} and returns a \texttt{ParseResult} wrapping an
    \texttt{ast::Program}, or reports parse errors and yields an empty result.
  \item \texttt{semantic::analyze} accepts the parsed program and produces an \texttt{AnalysisResult} containing the
    same \texttt{ast::Program} plus symbol-table and annotation side tables.
  \item \texttt{lowering::lower} consumes the \texttt{AnalysisResult} and returns a \texttt{LowerResult} holding an
    optional \texttt{ir::Program}.
  \item \texttt{midi::write\_midi} serializes the \texttt{ir::Program} to the output path; I/O failures are caught and
    reported as \texttt{DiagnosticStage::Output} errors.
\end{enumerate}

Figure~\ref{fig:compiler-pipeline-full} presents this control flow at full width, naming the C++ entry functions and
the guard conditions that terminate compilation early.

\subsection{Information Flow and Artifacts}\label{detail-subsec:information-flow-and-artifacts}

Each stage consumes a typed result object from its predecessor and produces the input expected by its successor.
Figure~\ref{fig:compiler-data-artifacts} summarizes the primary data structures and their fields; the annotated AST is
not a separate tree copy---semantic analysis attaches metadata to the existing \texttt{ast::Program} through
\texttt{AnalysisResult}.

\noindent \textbf{\texttt{ParseResult} (\texttt{motivo::parsing}).}
On success, owns a \texttt{std::unique\_ptr\textless ast::Program\textgreater{}} built by the Bison parser.
The program decomposes into a \texttt{Header} (tempo and time signature), a vector of \texttt{GlobalItem} nodes, and a
vector of \texttt{TrackDefinition} nodes.
The \texttt{ok()} predicate is true when the program pointer is non-null.

\noindent \textbf{\texttt{AnalysisResult} (\texttt{motivo::semantic}).}
Retains a const reference to the parsed \texttt{ast::Program} and owns private \texttt{SymbolTable} and
\texttt{Annotations} instances populated by \texttt{detail::Traversal::run()}.
Callers query resolved expression types and symbol identifiers through accessor methods such as
\texttt{get\_expression\_type} and \texttt{get\_resolved\_symbol}; lowering reads these caches instead of repeating
lookups.

\noindent \textbf{\texttt{ir::Program} (\texttt{motivo::lowering}).}
A flat timeline artifact: \texttt{tempo\_bpm}, \texttt{time\_sig\_numerator}, \texttt{time\_sig\_denominator}, and a
vector of \texttt{ir::Track} records.
Each track stores an \texttt{music::Instrument} and a chronologically sorted vector of \texttt{ir::NoteEvent} structures
(\texttt{midi\_note}, \texttt{start\_beat}, \texttt{duration\_beats}, \texttt{velocity}).

\noindent \textbf{MIDI output (\texttt{motivo::midi}).}
\texttt{write\_midi} does not return a structured result; it writes Standard MIDI File Format~1 bytes directly to disk
or throws an exception converted to an output-stage diagnostic.

\subsection{Transformation Stages}\label{detail-subsec:transformation-stages}

The transformation process is organized into four conceptual stages that map one-to-one onto the CMake library
boundaries described in Section~\ref{subsec:implementation-foundations}.

\noindent \textbf{Text to tree (frontend).} Flex (\texttt{parsing/detail/scanner.l}) tokenizes the character stream;
Bison (\texttt{parsing/detail/parser.y}) reduces tokens to an \texttt{ast::Program}.
The AST captures logical hierarchy (patterns, tracks, voices) without temporal resolution.

\noindent \textbf{Tree to annotated tree (semantic analysis).} \texttt{detail::Traversal} performs symbol resolution,
explicit type checking, and musical constraint validation in a single walk.
Collectors register pattern symbols before bodies are visited; validators enforce operand and call-site rules; visitors
in \texttt{semantic/detail/traversal/visitors/} dispatch on \texttt{std::variant} node kinds.

\noindent \textbf{Tree to timeline (lowering).} \texttt{detail::LowererContext} executes the annotated program like an
interpreter: \texttt{ast\_lowerers/} emit \texttt{NoteEvent}s, \texttt{expression\_evaluators/} resolve runtime
\texttt{ir::Value}s, and \texttt{value\_flattener.cpp} expands chords and sequences into atomic events.

\noindent \textbf{Timeline to binary (backend).} \texttt{midi/write\_midi.cpp} converts absolute beat positions to MIDI
ticks, splits each note into Note-On and Note-Off messages, stable-sorts events, and writes delta-time encoded track
chunks.

\begin{figure}[p]
  \centering
  \motivoscalefig{%
    \begin{tikzpicture}[
        stage/.style={motivo/stage, minimum width=10.4cm, minimum height=1.25cm},
        guard/.style={motivo/guard},
        arrow/.style={motivo/arrow}
      ]
      \node[stage, fill=motivoSoft] (input) {\textbf{Input stream}\\\texttt{FILE*} + source name + output path};
      \node[stage, fill=motivoBlue!16, below=0.7cm of input] (parse)
      {\textbf{Parsing}\\\texttt{parsing::parse\_stream}\\Flex \texttt{scanner.l} + Bison \texttt{parser.y}};
      \node[stage, fill=motivoCyan!20, below=0.7cm of parse] (sem) {\textbf{Semantic
      analysis}\\\texttt{semantic::analyze} $\rightarrow$ \texttt{detail::Traversal::run()}};
      \node[stage, fill=motivoGreen!20, below=0.7cm of sem] (low) {\textbf{Lowering}\\\texttt{lowering::lower}
      $\rightarrow$ \texttt{detail::LowererContext}};
      \node[stage, fill=motivoAmber!24, below=0.7cm of low] (midi) {\textbf{MIDI
      serialization}\\\texttt{midi::write\_midi}};
      \node[stage, fill=motivoSoft, below=0.7cm of midi] (out) {\textbf{Output artifact}\\Standard MIDI File
      \texttt{.mid}};

      \draw[arrow] (input) -- (parse);
      \draw[arrow] (parse) -- node[guard, left=0.12cm, align=right] {abort if \texttt{!parse\_result.ok()}} (sem);
      \draw[arrow] (sem) -- node[guard, left=0.12cm, align=right] {abort if semantic errors} (low);
      \draw[arrow] (low) -- node[guard, left=0.12cm, align=right] {abort if \texttt{!lowered.ok()}} (midi);
      \draw[arrow] (midi) -- (out);

      \node[guard, left=0.15cm of parse, align=right] {\texttt{DiagnosticsEngine}\\shared reference};
    \end{tikzpicture}%
  }
  \caption{Full compiler pipeline orchestrated by \texttt{motivo::compile} in \texttt{compiler.cpp},
  including early-exit guards after each stage.}\label{fig:compiler-pipeline-full}
\end{figure}

\begin{figure}[htbp]
  \centering
  \motivoscalefig{%
    \begin{tikzpicture}[
        box/.style={motivo/box, minimum width=2.85cm, minimum height=1.05cm},
        field/.style={motivo/guard},
        arrow/.style={motivo/arrow}
      ]
      \node[box, fill=motivoBlue!14] (pr) {\textbf{ParseResult}\\\texttt{unique\_ptr\textless Program\textgreater}};
      \node[box, fill=motivoCyan!18, right=0.45cm of pr] (ar) {\textbf{AnalysisResult}\\\texttt{Program
      \&}\\\texttt{SymbolTable}\\\texttt{Annotations}};
      \node[box, fill=motivoGreen!18, right=0.45cm of ar] (ir) {\textbf{ir::Program}\\\texttt{Header
      fields}\\\texttt{vector\textless Track\textgreater}\\\texttt{vector\textless NoteEvent\textgreater}};
      \node[box, fill=motivoAmber!22, right=0.45cm of ir] (mid) {\textbf{SMF bytes}\\\texttt{MThd + MTrk}\\480 PPQ};

      \draw[arrow] (pr) -- node[field, above] {analyze} (ar);
      \draw[arrow] (ar) -- node[field, above] {lower} (ir);
      \draw[arrow] (ir) -- node[field, above] {write\_midi} (mid);

      \node[field, below=0.4cm of pr] {\texttt{ast::Header}\\\texttt{globals[]}\\\texttt{tracks[]}};
      \node[field, below=0.4cm of ar] {same AST nodes\\+ side tables};
      \node[field, below=0.4cm of ir] {flat timeline\\no control flow};
      \node[field, below=0.4cm of mid] {relative delta-times\\per track};
    \end{tikzpicture}%
  }
  \caption{Stage artifacts passed between compiler modules: \texttt{ParseResult} $\rightarrow$
    \texttt{AnalysisResult} $\rightarrow$ \texttt{ir::Program} $\rightarrow$ MIDI
  file.}\label{fig:compiler-data-artifacts}
\end{figure}

\subsection{Diagnostics Flow}\label{detail-subsec:diagnostics-flow}

Diagnostics are centralized in \texttt{motivo\_diagnostics} (\texttt{common/diagnostics/diagnostics\_engine.cpp}).
Each reporting site calls \texttt{DiagnosticsEngine::report} with a stage, severity, optional \texttt{source::Location},
and message string.
Lexical errors originate in \texttt{scanner.l} via \texttt{report\_lexical}; syntax errors use Bison's
\texttt{parse.error
detailed} mode; semantic errors are emitted from \texttt{Traversal::diagnose}; lowering errors flow through
\texttt{LowererContext::report\_lowering\_error}; and file I/O failures in \texttt{write\_midi} become
\texttt{DiagnosticStage::Output} entries.
The final \texttt{CompileResult} returned to callers contains the flattened \texttt{Diagnostics} vector produced by
\texttt{diagnostics.take\_diagnostics()}.

\subsection{Technical Foundations and Build Structure}\label{subsec:implementation-foundations}

The implementation balances high-level musical expressiveness with low-level binary precision.
The discussion below covers the language and build choices that support this architecture.

\subsection{Language Choice: Modern C++ (C++23)}\label{detail-subsec:language-choice:-modern-c++-(c++23)}

The compiler targets the C++23 standard.
C++ was selected for deterministic performance during recursive unrolling, a strong static type system for AST and IR
modeling, and native support for binary serialization in the MIDI backend.

\subsubsection{C++23 Features Used in the Architecture}

\noindent \textbf{\texttt{std::variant} as a sum type.} AST and IR nodes use \texttt{std::variant} (e.g.,
\texttt{StatementKind}, \texttt{ExpressionKind}, \texttt{ir::Value}) with \texttt{std::visit} dispatch in semantic
visitors, lowering lowerers, and expression evaluators.

\noindent \textbf{Value-oriented memory management.} Recursive AST nodes use \texttt{std::unique\_ptr}; side tables
(\texttt{SymbolTable}, \texttt{Annotations}) use contiguous vectors indexed by opaque ID types
(\texttt{SymbolId}, \texttt{ScopeId}).

\noindent \textbf{Standard ranges.} \texttt{std::ranges::stable\_sort} orders MIDI events; \texttt{std::ranges::any\_of}
checks diagnostic severity in \texttt{CompileResult::has\_errors()}.

\subsection{Modular Build Infrastructure (CMake)}\label{detail-subsec:modular-build-infrastructure-(cmake)}

CMake~\cite{CMake_2026} builds the compiler as a graph of static libraries declared under
\texttt{compiler/src/CMakeLists.txt} and \texttt{compiler/src/motivo/*/CMakeLists.txt}.
Flex and Bison generate \texttt{scanner.cpp} and \texttt{parser.cpp} into the build tree; the scanner includes the
generated \texttt{parser.hpp} header, so Bison runs before Flex.

\subsubsection{Library Decomposition}

Figure~\ref{fig:compiler-module-graph} shows link dependencies as declared in CMake.
The names below match the \texttt{add\_library} targets in the repository.

\noindent \textbf{\texttt{motivo\_location} (\texttt{common/source/}).} Source positions (\texttt{Location},
\texttt{Position}) attached to AST nodes and diagnostics.

\noindent \textbf{\texttt{motivo\_diagnostics} (\texttt{common/diagnostics/}).} Public
\texttt{DiagnosticsEngine}; depends
on \texttt{motivo\_location}.
Linked by every stage library and exposed through \texttt{include/motivo/diagnostics/diagnostic.hpp}.

\noindent \textbf{\texttt{motivo\_music} (\texttt{common/music/}).} Pitch, accidental, note literals, instruments, and
drum-note enums shared by parsing, lowering, and MIDI.

\noindent \textbf{\texttt{motivo\_types} (\texttt{common/types/}).} The \texttt{types::Type} enum and
\texttt{type\_rules} module (\texttt{is\_assignable}, \texttt{numeric\_result}, \texttt{binary\_result\_type}).

\noindent \textbf{\texttt{motivo\_parsing} (\texttt{parsing/}).} Flex/Bison generated sources, \texttt{parse.cpp},
\texttt{parse\_result.cpp}, and scanner helpers.
Public dependency: \texttt{motivo\_diagnostics}; private: \texttt{motivo\_location}, \texttt{motivo\_music}.
AST node definitions live in headers under \texttt{common/ast/} compiled into downstream modules.

\noindent \textbf{\texttt{motivo\_semantic} (\texttt{semantic/}).} \texttt{analyze.cpp}, \texttt{analysis\_result.cpp},
\texttt{traversal.cpp}, collectors, validators, and visitor translation units under \texttt{detail/}.
Depends on \texttt{motivo\_diagnostics}, \texttt{motivo\_location}, and \texttt{motivo\_types}.

\noindent \textbf{\texttt{motivo\_lowering} (\texttt{lowering/}).} \texttt{lower.cpp}, \texttt{lowerer\_context.cpp},
\texttt{value\_flattener.cpp}, \texttt{expression\_evaluator.cpp}, and split lowerer/evaluator files under
\texttt{detail/}.
Depends on \texttt{motivo\_semantic} (for \texttt{AnalysisResult}) and \texttt{motivo\_music}.

\noindent \textbf{\texttt{motivo\_midi} (\texttt{midi/}).} \texttt{write\_midi.cpp} only; depends on
\texttt{motivo\_lowering} for \texttt{ir::Program} and instrument enums.

\noindent \textbf{\texttt{motivo\_compiler} (\texttt{compiler.cpp}).} Top-level orchestrator linking all stage
libraries; consumed by the \texttt{motivoc} CLI application.

\subsubsection{Directory Organization and Public API}

\noindent \textbf{\texttt{include/motivo/}.} Minimal public surface: \texttt{compiler.hpp} and
\texttt{diagnostics/diagnostic.hpp}.

\noindent \textbf{\texttt{src/motivo/}.} All stage implementations, generated parser/scanner sources (build tree), and
internal headers (e.g., \texttt{semantic/analysis\_result.hpp}, \texttt{lowering/lower.hpp}).

The \texttt{motivo::compile} function hides this modular complexity behind a single call suitable for the CLI and the
Motivo Studio compile API.

\begin{figure}[htbp]
  \centering
  \motivoscalefig{%
    \begin{tikzpicture}[
        mod/.style={motivo/mod, minimum width=2.45cm, minimum height=0.78cm},
        root/.style={motivo/mod, minimum width=2.95cm, font=\small\sffamily},
        dep/.style={-Stealth, thick, draw=motivoInk!55}
      ]
      \node[root, fill=motivoSoft] (comp) {\texttt{motivo\_compiler}};
      \node[mod, fill=motivoBlue!14, below left=0.55cm and 0.15cm of comp] (parse) {\texttt{motivo\_parsing}};
      \node[mod, fill=motivoCyan!18, below=0.55cm of comp] (sem) {\texttt{motivo\_semantic}};
      \node[mod, fill=motivoGreen!18, below right=0.55cm and 0.15cm of comp] (low) {\texttt{motivo\_lowering}};
      \node[mod, fill=motivoAmber!22, right=0.45cm of low] (midi) {\texttt{motivo\_midi}};

      \node[mod, fill=motivoSoft, below=0.55cm of parse] (diag) {\texttt{motivo\_diagnostics}};
      \node[mod, fill=motivoSoft, left=0.35cm of diag] (loc) {\texttt{motivo\_location}};
      \node[mod, fill=motivoSoft, right=0.35cm of diag] (music) {\texttt{motivo\_music}};
      \node[mod, fill=motivoSoft, below=0.55cm of sem] (types) {\texttt{motivo\_types}};

      \draw[dep] (comp) -- (parse);
      \draw[dep] (comp) -- (sem);
      \draw[dep] (comp) -- (low);
      \draw[dep] (comp) -- (midi);
      \draw[dep] (parse) -- (diag);
      \draw[dep] (sem) -- (diag);
      \draw[dep] (low) -- (diag);
      \draw[dep] (diag) -- (loc);
      \draw[dep] (parse) -- (music);
      \draw[dep] (low) -- (music);
      \draw[dep] (midi) -- (low);
      \draw[dep] (sem) -- (types);
      \draw[dep] (low) -- (sem);
    \end{tikzpicture}%
  }
  \caption{CMake static-library dependency graph for the Motivo compiler (primary
  \texttt{target\_link\_libraries} edges).}\label{fig:compiler-module-graph}
\end{figure}
